\chapter{Papildomi uždaviniai}

Šiame skyriuje pateikiami papildomi VBE stiliaus uždaviniai visiems pagrindiniams
12 klasės matematikos kursui temoms. Kiekvienas uždavinys atitinka egzamino
formatą ir sudėtingumo lygį.

% ============================================================
\section*{A blokas — Trumpieji uždaviniai (po 1 tašką)}
% ============================================================

\begin{exercise}
Apskaičiuokite:
\[
\log_3 \frac{1}{27} + \log_3 9.
\]
\end{exercise}

\begin{exercise}
Suprastinkite išraišką ($a > 0$):
\[
\frac{a^{3/2} \cdot a^{-1/2}}{a^{1/4} \cdot a^{3/4}}.
\]
\end{exercise}

\begin{exercise}
Išspręskite lygtį:
\[
2^{2x-1} = 8.
\]
\end{exercise}

\begin{exercise}
Nustatykite funkcijos $f(x) = \sqrt{5 - 2x}$ apibrėžimo sritį.
\end{exercise}

\begin{exercise}
Apskaičiuokite $\sin 150^\circ \cdot \cos 60^\circ - \cos 150^\circ \cdot \sin 60^\circ$.
\end{exercise}

\begin{exercise}
Raskite aritmetinės progresijos $(a_n)$ pirmą narį, jei $a_5 = 17$ ir $d = 3$.
\end{exercise}

\begin{exercise}
Geometrinės progresijos $(b_n)$ pirmasis narys $b_1 = 6$, vardiklis $q = \tfrac{1}{2}$.
Apskaičiuokite begalinės geometrinės progresijos sumą $S$.
\end{exercise}

\begin{exercise}
Išspręskite nelygybę:
\[
x^2 - 5x + 6 < 0.
\]
\end{exercise}

\begin{exercise}
Trikampyje $ABC$ žinoma, kad $AB = 5$, $BC = 7$, $\angle ABC = 60^\circ$.
Apskaičiuokite trikampio plotą.
\end{exercise}

\begin{exercise}
Apskaičiuokite išvestinę $f'(x)$, jei $f(x) = 3x^4 - 2x^2 + 5x - 1$.
\end{exercise}

\begin{exercise}
Apskaičiuokite integralą:
\[
\int_0^2 (3x^2 - 2x + 1)\,dx.
\]
\end{exercise}

\begin{exercise}
Aibė $A = \{1,\, 2,\, 3,\, 4,\, 5\}$ ir $B = \{3,\, 4,\, 5,\, 6,\, 7\}$.
Raskite $A \cap B$ ir $A \cup B$.
\end{exercise}

\begin{exercise}
Kauliukas metamas vieną kartą. Kokia tikimybė, kad iškris lyginis skaičius?
\end{exercise}

\begin{exercise}
Duota, kad $\tg \alpha = 2$ ir $\alpha \in \left(0;\, \tfrac{\pi}{2}\right)$.
Apskaičiuokite $\sin \alpha$ ir $\cos \alpha$.
\end{exercise}

\begin{exercise}
Raskite skaičių $n$, jei $C_n^2 = 15$.
\end{exercise}

\begin{exercise}
Išskleiskite skliaustus ir suprastinkite:
\[
(x + 3)^2 - (x - 3)^2.
\]
\end{exercise}

\begin{exercise}
Nustatykite, kuriam $k$ reikšmių rinkiniui tiesė $y = kx + 2$ nesikirčia su
parabole $y = x^2 + 3x + 4$.
\end{exercise}

\begin{exercise}
Apskaičiuokite $\lim_{x \to 2} \dfrac{x^2 - 4}{x - 2}$.
\end{exercise}

\begin{exercise}
Apskaičiuokite kūgio tūrį, jei jo pagrindo spindulys $r = 3$\,cm ir aukštis
$h = 4$\,cm.
\end{exercise}

\begin{exercise}
Duota funkcija $f(x) = \ln(x^2 + 1)$. Raskite $f'(x)$.
\end{exercise}


% ============================================================
\section*{B blokas — Daugiadaliai uždaviniai (po 6–8 taškus)}
% ============================================================

% ------- 1. Lygtys ir nelygybės -------
\begin{exercise}
\textbf{Lygtys ir nelygybės.}
\begin{enumerate}[label=\arabic*.]
  \item Išspręskite lygtį:
  \[
    \log_2(x+3) + \log_2(x-1) = 3.
  \]
  \item Išspręskite nelygybę:
  \[
    \frac{x^2 - x - 6}{x + 4} \geq 0.
  \]
  \item Išspręskite lygtį:
  \[
    2\sin^2 x - 3\sin x + 1 = 0, \quad x \in [0;\, 2\pi].
  \]
\end{enumerate}
\end{exercise}

% ------- 2. Funkcijų tyrimas -------
\begin{exercise}
\textbf{Funkcijų tyrimas ir grafikai.}\\
Duota funkcija
\[
  f(x) = x^3 - 3x^2 - 9x + 2.
\]
\begin{enumerate}[label=\arabic*.]
  \item Raskite funkcijos $f$ išvestinę $f'(x)$.
  \item Nustatykite, kuriuose taškuose funkcija turi lokalius ekstremumus.
        Apskaičiuokite jų reikšmes.
  \item Nustatykite funkcijos didėjimo ir mažėjimo intervalus.
  \item Nupieškite apytikslį funkcijos grafiką intervale $[-2;\, 5]$.
\end{enumerate}
\end{exercise}

% ------- 3. Integralai ir plotai -------
\begin{exercise}
\textbf{Integralas ir figūros plotas.}\\
Duotos dvi funkcijos: $f(x) = x^2$ ir $g(x) = 2x + 3$.
\begin{enumerate}[label=\arabic*.]
  \item Raskite funkcijų grafiko susikirtimo taškus.
  \item Apskaičiuokite figūros, apribotos abiem kreivėmis, plotą naudodami
        integralą.
  \item Apskaičiuokite $\displaystyle\int_{-1}^{3} \bigl(g(x) - f(x)\bigr)\,dx$ ir
        palyginkite su gautu plotu.
\end{enumerate}
\end{exercise}

% ------- 4. Eksponentinė ir logaritminė funkcija -------
\begin{exercise}
\textbf{Eksponentinė ir logaritminė funkcija.}
\begin{enumerate}[label=\arabic*.]
  \item Išspręskite lygtį:
  \[
    3^{x+1} + 3^{x-1} = 30.
  \]
  \item Išspręskite nelygybę:
  \[
    \ln(2x - 1) < \ln(x + 4).
  \]
  \item Duota, kad $\log_a 2 = m$ ir $\log_a 3 = n$. Išreikškite $\log_a 72$
        per $m$ ir $n$.
\end{enumerate}
\end{exercise}

% ------- 5. Trigonometrija -------
\begin{exercise}
\textbf{Trigonometriniai uždaviniai.}
\begin{enumerate}[label=\arabic*.]
  \item Įrodykite tapatybę:
  \[
    \frac{\sin x}{1 + \cos x} + \frac{1 + \cos x}{\sin x} = \frac{2}{\sin x}.
  \]
  \item Išspręskite lygtį:
  \[
    \cos 2x - 3\cos x + 2 = 0.
  \]
  \item Apskaičiuokite $\sin\!\left(\alpha + \beta\right)$, jei
        $\sin\alpha = \tfrac{3}{5}$, $\cos\beta = \tfrac{5}{13}$,
        $\alpha, \beta \in \left(0;\,\tfrac{\pi}{2}\right)$.
\end{enumerate}
\end{exercise}

% ------- 6. Sekos ir eilutės -------
\begin{exercise}
\textbf{Sekos.}\\
Aritmetinės progresijos pirmasis narys $a_1 = 4$, o jos suma $S_{10} = 175$.
\begin{enumerate}[label=\arabic*.]
  \item Raskite progresijos skirtumą $d$.
  \item Apskaičiuokite $a_{20}$.
  \item Raskite, koks yra mažiausias $n$, kuriam $a_n > 100$.
\end{enumerate}
\end{exercise}

% ------- 7. Kombinatorika ir tikimybė -------
\begin{exercise}
\textbf{Kombinatorika ir tikimybė.}
\begin{enumerate}[label=\arabic*.]
  \item Iš 5 berniukų ir 4 mergaičių reikia sudaryti 4 asmenų komandą.
        Kiek yra galimų komandų, jeigu komandoje turi būti lygiai 2
        berniukai ir 2 mergaitės?
  \item Atsitiktinai pasirenkamas skaičius iš aibės $\{1,\, 2,\, 3,\, \ldots,\, 20\}$.
        Kokia tikimybė, kad pasirinktas skaičius dalijasi iš 3 arba iš 4?
  \item Moneta metama 4 kartus. Apskaičiuokite tikimybę, kad herbas iškris
        lygiai 3 kartus.
\end{enumerate}
\end{exercise}

% ------- 8. Planimetrija -------
\begin{exercise}
\textbf{Planimetrija.}\\
Duotas trikampis $ABC$, kuriame $AB = 8$, $AC = 6$, $\angle BAC = 60^\circ$.
\begin{enumerate}[label=\arabic*.]
  \item Raskite kraštinę $BC$ naudodami kosinusų teoremą.
  \item Apskaičiuokite trikampio plotą.
  \item Apskaičiuokite į trikampį įbrėžto apskritimo spindulį.
\end{enumerate}
\end{exercise}

% ------- 9. Stereometrija -------
\begin{exercise}
\textbf{Stereometrija.}\\
Taisyklingosios trikampės piramidės pagrindo kraštinė $a = 6$, o šoninė
kraštinė $l = 5$.
\begin{enumerate}[label=\arabic*.]
  \item Apskaičiuokite piramidės aukštį $h$.
  \item Apskaičiuokite piramidės tūrį $V$.
  \item Apskaičiuokite pilnąjį paviršiaus plotą $S_{\text{pil}}$.
\end{enumerate}
\end{exercise}

% ------- 10. Optimizavimas -------
\begin{exercise}
\textbf{Optimizavimo uždavinys.}\\
Stačiakampio perimetras lygus $40$\,cm.
\begin{enumerate}[label=\arabic*.]
  \item Išreikškite stačiakampio plotą $S$ kaip vienos kraštinės $x$ funkciją.
  \item Nustatykite, kuriam $x$ plotas yra didžiausias.
  \item Apskaičiuokite didžiausią plotą.
\end{enumerate}
\end{exercise}

% ------- 11. Tiesinė algebra ir koordinatės -------
\begin{exercise}
\textbf{Koordinačių geometrija.}\\
Duotos dvi tiesės: $\ell_1\colon 2x - y + 1 = 0$ ir $\ell_2\colon x + 3y - 7 = 0$.
\begin{enumerate}[label=\arabic*.]
  \item Raskite tiesių susikirtimo tašką.
  \item Nustatykite kiekvienos tiesės nuolydį ($k$ koeficientą).
  \item Apskaičiuokite kampą tarp tiesių $\ell_1$ ir $\ell_2$.
\end{enumerate}
\end{exercise}

% ------- 12. Įrodymas (aukštesnio lygio) -------
\begin{exercise}
\textbf{Matematinis įrodymas.}
\begin{enumerate}[label=\arabic*.]
  \item Įrodykite, kad bet kurių trijų nuoseklių lyginių skaičių sandauga
        dalijasi iš $48$.
  \item Tarkime, kad $a + b + c = 0$. Įrodykite, jog
  \[
    a^3 + b^3 + c^3 = 3abc.
  \]
  \item Naudodami matematinę indukciją, įrodykite, kad visoms natūralioms
        $n$ reikšmėms:
  \[
    1 + 3 + 5 + \cdots + (2n-1) = n^2.
  \]
\end{enumerate}
\end{exercise}

% ------- 13. Sudėtingas kompleksinis uždavinys -------
\begin{exercise}
\textbf{Kompleksinis uždavinys.}\\
Vienas darbininkas gali baigti darbą per $x$ valandų, kitas — per $(x + 2)$
valandas.
\begin{enumerate}[label=\arabic*.]
  \item Išreikškite per 1 valandą kiekvieno darbininko atliktą darbo dalį.
  \item Sudarykite lygtį, jei dirbdami kartu jie darbą atlieka per $\frac{15}{4}$
        valandas.
  \item Išspręskite lygtį ir raskite, per kiek valandų kiekvienas darbininkas
        atliktų darbą vienas.
\end{enumerate}
\end{exercise}

% ------- 14. Skaičių teorija -------
\begin{exercise}
\textbf{Skaičių teorija ir skaičiavimai.}\\
Duota: $7^n - 5^n = K$ ir $7^n + 5^n = L$.
\begin{enumerate}[label=\arabic*.]
  \item Išreikškite $7^{2n} - 5^{2n}$ per $K$ ir $L$.
  \item Raskite $35^n$, jei $K = 2$ ir $L = 10$.
  \item Patikrinkite: ar $K^2 + L^2$ visada lyginė? Pagrįskite.
\end{enumerate}
\end{exercise}

% ------- 15. Statistika ir tikimybių pasiskirstymas -------
\begin{exercise}
\textbf{Tikimybių pasiskirstymas ir statistika.}\\
Mokytoja žino, kad jos klasėje kiekvienas mokinys gali atsakyti testo
klausimą teisingai su tikimybe $p = 0{,}7$. Klausimyne yra $10$ klausimų.
\begin{enumerate}[label=\arabic*.]
  \item Apskaičiuokite tikimybę, kad mokinys atsakys lygiai $7$ iš $10$
        klausimų teisingai (naudokite Bernulio formulę).
  \item Apskaičiuokite teisingų atsakymų skaičiaus vidurkį (tikėtinąją
        reikšmę) $E(X)$.
  \item Apskaičiuokite dispersiją $D(X)$ ir standartinį nuokrypį $\sigma$.
\end{enumerate}
\end{exercise}

% ------- 16. Funkcijų analizė — sudėtingesnė -------
\begin{exercise}
\textbf{Sudėtingesnė funkcijų analizė.}\\
Duota funkcija $f(x) = xe^{-x}$.
\begin{enumerate}[label=\arabic*.]
  \item Raskite $f'(x)$ ir $f''(x)$.
  \item Nustatykite ekstremumus ir infleksijos taškus.
  \item Nustatykite, ar funkcija turi asimptotą, kai $x \to +\infty$.
        Argumentuokite atsakymą.
  \item Nupieškite apytikslį funkcijos grafiką.
\end{enumerate}
\end{exercise}

% ------- 17. Vektoriai ir koordinatės -------
\begin{exercise}
\textbf{Vektoriai ir tiesė erdvėje.}\\
Duoti taškai $A(1;\,2;\,3)$, $B(4;\,0;\,-1)$ ir $C(-1;\,3;\,2)$.
\begin{enumerate}[label=\arabic*.]
  \item Raskite vektorius $\overrightarrow{AB}$ ir $\overrightarrow{AC}$.
  \item Apskaičiuokite vektorių skaliarinę sandaugą
        $\overrightarrow{AB} \cdot \overrightarrow{AC}$.
  \item Apskaičiuokite kampą tarp vektorių $\overrightarrow{AB}$ ir
        $\overrightarrow{AC}$.
  \item Raskite trikampio $ABC$ plotą.
\end{enumerate}
\end{exercise}

% ------- 18. Sfera ir kūgis -------
\begin{exercise}
\textbf{Rutulys ir kūgis.}\\
Rutulys yra įbrėžtas į tiesųjį kūgį. Kūgio pagrindo spindulys $R = 6$\,cm,
o aukštis $H = 8$\,cm.
\begin{enumerate}[label=\arabic*.]
  \item Apskaičiuokite kūgio šoninę kraštinę $l$.
  \item Raskite įbrėžto rutulio spindulį $r$.
  \item Apskaičiuokite rutulio ir kūgio tūrių santykį $\dfrac{V_r}{V_k}$.
\end{enumerate}
\end{exercise}

% ------- 19. Parametrinis uždavinys -------
\begin{exercise}
\textbf{Lygtys su parametru.}\\
Duota kvadratinė lygybė $x^2 - 2ax + (a^2 - a - 2) = 0$.
\begin{enumerate}[label=\arabic*.]
  \item Apskaičiuokite diskriminantą $\Delta$ ir išreikškite per $a$.
  \item Raskite visas $a$ reikšmes, kurioms lygybė turi du skirtingus
        realius šaknis.
  \item Raskite $a$ reikšmę, kuriai abu šaknys yra teigiamos.
\end{enumerate}
\end{exercise}

% ------- 20. Pilnas sprendimas — aukštesnio lygio -------
\begin{exercise}
\textbf{Kompleksinis aukšto lygio uždavinys.}\\
Ūkininkas turi $120$\,m tvoros ir nori apitverti stačiakampį daržą, kurio
viena pusė yra palei upę (tos pusės tverti nereikia).
\begin{enumerate}[label=\arabic*.]
  \item Nustatykite kintamąjį ir sudarykite ploto funkciją.
  \item Apskaičiuokite ploto funkciją naudodami diferencialinį skaičiavimą
        ir raskite, kokių matmenų darzas turi būti, kad jo plotas būtų
        didžiausias.
  \item Apskaičiuokite didžiausią galimą plotą.
  \item Įrodykite, kad rastas taškas tikrai yra maksimumas (ne minimumas),
        naudodami antrąją išvestinę.
\end{enumerate}
\end{exercise}

